\documentclass[UTF8,fontset=none,11pt]{ctexart}
\usepackage[a4paper,margin=2.05cm]{geometry}
\usepackage{amsmath,amssymb}
\usepackage{enumitem}
\usepackage[table]{xcolor}
\usepackage{booktabs}
\usepackage{array}
\usepackage{longtable}
\usepackage{tikz}
\usepackage{graphicx}
\usepackage{hyperref}
\usetikzlibrary{positioning,arrows.meta,shapes.geometric,fit,backgrounds,calc}

\setCJKmainfont[AutoFakeSlant=0.18]{PingFang SC}
\setCJKsansfont[AutoFakeSlant=0.18]{PingFang SC}
\setCJKmonofont[AutoFakeSlant=0.18]{PingFang SC}

\hypersetup{
  colorlinks=true,
  linkcolor=blue!55!black,
  urlcolor=blue!55!black,
  citecolor=blue!55!black,
  pdftitle={Semantic Analysis Intro SDD SDT 中文学习笔记},
  pdfauthor={Trae AI}
}

\setlength{\parindent}{0pt}
\setlength{\parskip}{0.55em}
\setlist[itemize]{leftmargin=1.8em,itemsep=0.22em,topsep=0.25em}
\setlist[enumerate]{leftmargin=2.0em,itemsep=0.28em,topsep=0.3em}
\renewcommand{\arraystretch}{1.18}
\setlength{\emergencystretch}{3em}
\sloppy

\definecolor{defrule}{RGB}{25,92,140}
\definecolor{noterule}{RGB}{120,80,15}
\definecolor{warnrule}{RGB}{160,50,45}
\definecolor{softblue}{RGB}{235,246,255}
\definecolor{softyellow}{RGB}{255,249,230}
\definecolor{graytext}{RGB}{95,95,95}
\definecolor{GRAYTEXT}{RGB}{95,95,95}
\newcolumntype{L}[1]{>{\raggedright\arraybackslash}p{#1}}

\newenvironment{defbox}[1][]{%
  \par\medskip\noindent{\color{defrule}\rule{\linewidth}{1.1pt}}\par\nopagebreak
  \noindent{\bfseries\color{defrule}#1}\par\nopagebreak\smallskip}%
  {\par\nopagebreak\smallskip\noindent{\color{defrule}\rule{\linewidth}{0.4pt}}\par\medskip}
\newenvironment{notebox}[1][]{%
  \par\medskip\noindent{\color{noterule}\rule{\linewidth}{1.1pt}}\par\nopagebreak
  \noindent{\bfseries\color{noterule}#1}\par\nopagebreak\smallskip}%
  {\par\nopagebreak\smallskip\noindent{\color{noterule}\rule{\linewidth}{0.4pt}}\par\medskip}
\newenvironment{warnbox}[1][]{%
  \par\medskip\noindent{\color{warnrule}\rule{\linewidth}{1.1pt}}\par\nopagebreak
  \noindent{\bfseries\color{warnrule}#1}\par\nopagebreak\smallskip}%
  {\par\nopagebreak\smallskip\noindent{\color{warnrule}\rule{\linewidth}{0.4pt}}\par\medskip}

\newcommand{\pg}[1]{\texorpdfstring{{\small\color{graytext}\,[p#1]}}{ [p#1]}}
\newcommand{\eps}{\varepsilon}
\newcommand{\FIRST}{\ensuremath{\mathrm{FIRST}}}
\newcommand{\FOLLOW}{\ensuremath{\mathrm{FOLLOW}}}
\newcommand{\SDD}{\ensuremath{\mathrm{SDD}}}
\newcommand{\SDT}{\ensuremath{\mathrm{SDT}}}
\newcommand{\LR}{\ensuremath{\mathrm{LR}}}
\newcommand{\LL}{\ensuremath{\mathrm{LL}}}
\newcommand{\SLR}{\ensuremath{\mathrm{SLR}}}
\newcommand{\ACTION}{\ensuremath{\mathrm{ACTION}}}
\newcommand{\GOTO}{\ensuremath{\mathrm{GOTO}}}

\title{\bfseries Lecture 6: Semantic Analysis Intro, SDD \& SDT\\中文学习笔记}
\author{基于课件 \texttt{6.Semantic\_Analysis\_Intro\_SDD\_SDT.pdf} 整理}
\date{\today}

\begin{document}
\maketitle
\tableofcontents
\newpage

\section*{阅读导引}
\addcontentsline{toc}{section}{阅读导引}

本讲从语义分析进入编译前端的最后阶段。前几讲的词法分析和语法分析主要回答“源程序长得合不合法”；本讲开始回答“它的含义是否合理、如何把语义信息带给后续阶段”。核心工具是语法制导定义（SDD）和语法制导翻译方案（SDT），再进一步讨论如何在 LL/LR 分析器中实现属性计算，以及语义分析的典型应用：绑定、作用域、符号表和类型检查。

\begin{notebox}[本讲主线]
语义分析 = 给语法树上的符号补充语义属性，并检查/传播这些属性。SDD 描述“属性怎么算”，SDT 描述“动作什么时候执行”。若要在语法分析过程中完成语义分析，就必须让属性依赖顺序与分析器的执行顺序兼容。
\end{notebox}

\section{编译阶段与语义分析位置 \pg{1--4}}

语义分析（semantic analysis）位于编译前端的最后阶段。整体流程为：
\[
\text{字符流}\to \text{词法分析}\to \text{token stream}\to \text{语法分析}\to \text{syntax tree}\to \text{语义分析}.
\]
语义分析之后，编译器进入中间代码生成、优化和目标代码生成。

课件展示了 AST 上带类型等标注的例子。普通 AST 只表达结构，例如 while、赋值、加法、比较等节点；Annotated AST 或 Decorated AST 会在节点上额外标注类型、符号表入口、地址、作用域等信息。后续中间代码生成通常依赖这些语义信息。

\begin{defbox}[Annotated AST / Decorated AST]
带标注的 AST 是在抽象语法树节点上附加语义属性后的树。它不只是告诉我们程序结构，还告诉我们每个标识符绑定到哪个定义、表达式类型是什么、变量存储在哪里等。
\end{defbox}

\section{为什么需要语义分析 \pg{5--8}}

\subsection{上下文决定含义 \pg{5--6}}

程序使用大量符号，也就是标识符（identifier）。单看一个标识符的名字，编译器无法知道它表示哪个变量、函数、字段或类型；必须结合上下文判断。

课件用英文句子类比：``He ate it'' 在语法上正确，但只有知道前文是 ``Sam bought a pizza'' 时，``He'' 才能指 Sam，``it'' 才能指 pizza；如果前文是 ``Sam bought a car''，``it'' 就不能被 ate 合理使用。

语义分析类似地做两件事：

\begin{itemize}
  \item 关联：把标识符使用点和它引用的对象联系起来，例如把变量使用绑定到某个声明。
  \item 检查：判断使用是否合法，例如变量是否先声明后使用，类型是否支持对应操作。
\end{itemize}

语义比语法更难描述。语法只描述程序形式是否符合文法；语义描述程序执行时的意义。CFG 无法处理很多上下文关系，例如“某个使用点必须匹配前面某个同名定义”。课件指出，这类 def-use 模式类似 \(\{wcw\mid w\in(a|b)^*\}\)，不是普通 CFG parser 能直接检查的内容。

\subsection{语义分析要做什么 \pg{7}}

语义分析是编译器前端最后一次拒绝错误程序的机会。它检查前面阶段无法捕获的性质，包括：

\begin{itemize}
  \item 变量必须先声明再使用。
  \item 标识符使用时类型一致。
  \item 表达式类型正确，例如布尔表达式不能当整数加法操作数。
  \item 函数调用参数个数和类型匹配。
  \item 继承层次、作用域中变量数量等后续阶段需要的信息。
\end{itemize}

\subsection{两类实现方式 \pg{8}}

\begin{longtable}{L{0.22\linewidth}L{0.34\linewidth}L{0.34\linewidth}}
\toprule
方式 & 思路 & 特点 \\
\midrule
属性文法 / Attribute grammar & 在语法分析过程中，根据产生式规则计算属性并做检查 & 可一遍完成，适合和 parser 紧密结合 \\
AST walk & 第一遍构造 parse tree 或 AST，第二遍遍历树做语义检查 & 阶段分离清晰，属性计算顺序更自由 \\
\bottomrule
\end{longtable}

如果先构造树再遍历，只要属性依赖没有环，理论上可以选择合适的遍历顺序；如果想在 parsing 过程中完成，就必须受 LL 或 LR 分析顺序限制。

\section{语法制导翻译、属性、SDD 与 SDT \pg{9--17}}

\subsection{语法制导翻译 \pg{9--11}}

\begin{defbox}[Syntax Directed Translation]
语法制导翻译（syntax directed translation）是以程序语法结构为驱动进行语义分析和翻译的方法。解析过程和 parse tree/AST 用来指导属性计算、语义检查和代码生成。
\end{defbox}

它的基本做法是增强原来的 CFG：

\begin{itemize}
  \item 给每个文法符号附加语义属性（semantic attributes）。
  \item 给每条产生式附加语义规则（semantic rules）或语义动作（semantic actions）。
\end{itemize}

这就是 CFG-driven translation：语法结构决定在哪里计算哪些语义信息。

\subsection{语义属性 \pg{12}}

\begin{defbox}[Attribute]
属性是附着在文法符号上的语义信息，记作 \(X.a\)，其中 \(X\) 是文法符号，\(a\) 是属性名。属性值可以是数字、字符串、类型、内存位置、寄存器、AST 指针、代码片段等。
\end{defbox}

常见用途：

\begin{itemize}
  \item 计算表达式值：属性可保存数值。
  \item 构造 AST：属性可保存指向 AST 节点的指针。
  \item 生成代码：属性可保存代码片段或临时变量名。
  \item 类型检查：属性可保存表达式或标识符的类型。
\end{itemize}

\subsection{SDD：语法制导定义 \pg{13--14}}

\begin{defbox}[Syntax Directed Definition]
SDD 是一个 CFG 加上文法符号的属性和产生式的语义规则。它是 CFG 的推广：CFG 只描述语法结构，SDD 还描述每个结构上的语义属性如何计算。
\end{defbox}

例子：
\[
E\to E_1+E_2,\qquad E.val=E_1.val+E_2.val.
\]
这里 \(val\) 是属性，语义规则说明父节点 \(E\) 的值由两个子表达式的值相加得到。

SDD 只描述依赖关系和计算公式，不要求规则在产生式右部的哪个位置执行。

\subsection{SDT：语法制导翻译方案 \pg{15}}

\begin{defbox}[Syntax Directed Translation Scheme]
SDT 是在 CFG 产生式右部嵌入程序片段的翻译方案，这些程序片段称为语义动作。语义动作在产生式右部的位置决定其执行时机。
\end{defbox}

例子：
\[
E\to E_1+E_2\ \{E.val=E_1.val+E_2.val\}.
\]
若动作在右部末尾，就表示整条右部都已识别后执行；若动作嵌在中间，就表示识别到该位置时立即执行。

\subsection{SDD 与 SDT 的区别 \pg{16--17}}

\begin{longtable}{L{0.18\linewidth}L{0.36\linewidth}L{0.36\linewidth}}
\toprule
对比项 & SDD & SDT \\
\midrule
本质 & CFG + 属性 + 语义规则 & CFG + 属性 + 嵌入在产生式中的语义动作 \\
定位 & 高层规则说明 & 可执行的翻译方案 \\
执行顺序 & 不直接规定，依赖属性依赖关系 & 由动作在 RHS 中的位置规定 \\
典型形式 & \(L.inh=T.type\) & \(D\to T\ \{L.inh=T.type\}\ L\) \\
\bottomrule
\end{longtable}

语义规则不绑定产生式右部的位置；语义动作必须在具体位置执行。若写成
\[
A\to \alpha\ \{action_1\}\ \beta\ \{action_2\}\ \gamma,
\]
则 \(action_1\) 在 \(\alpha\) 已产生、\(\beta\) 尚未处理前执行；\(action_2\) 在 \(\alpha, action_1, \beta\) 之后、\(\gamma\) 之前执行。

\section{SDD 属性类型与基本概念 \pg{18--28}}

\subsection{综合属性与继承属性 \pg{19--23}}

\begin{defbox}[综合属性]
综合属性（synthesized attribute）由某节点自身和它的子节点属性计算得到，信息通常从子节点向父节点上传。
\end{defbox}

例子：
\[
E\to E_1+T,\qquad E.val=E_1.val+T.val.
\]
父节点 \(E.val\) 依赖子节点 \(E_1.val\) 和 \(T.val\)，所以 \(val\) 是综合属性。

\begin{defbox}[继承属性]
继承属性（inherited attribute）由某节点的父节点、自身或兄弟节点属性计算得到，信息通常从父节点向子节点传递，或从左兄弟向右兄弟传递。
\end{defbox}

例子：
\[
D\to T L,\qquad L.inh=T.type.
\]
声明列表 \(L\) 继承类型 \(T.type\)，从而知道后续变量应该登记为 int 还是 float。

终结符可以有综合属性，例如词法分析器给 \(digit.lexval\)、\(id.lexeme\) 赋值；终结符没有继承属性，因为它们通常是词法分析输入提供的叶子，不再向下展开。

\subsection{例子：表达式值与变量类型 \pg{24--25}}

表达式求值 SDD：
\[
\begin{aligned}
E&\to E_1+T & E.val&=E_1.val+T.val,\\
T&\to T_1*F & T.val&=T_1.val*F.val,\\
F&\to digit & F.val&=digit.val.
\end{aligned}
\]
所有 \(val\) 都从叶子往上传，是综合属性。

变量声明 SDD：
\[
decl\to type\ varlist,\qquad varlist.dtype=type.dtype.
\]
类型信息从 \(type\) 传给变量列表，再传给每个 id，是继承属性。这样每个变量声明点都能被加入符号表，记录对应类型。

\subsection{副作用、属性文法和标注分析树 \pg{26--28}}

\begin{defbox}[Side effect]
副作用指属性值计算之外的额外行为，例如生成代码、打印结果、修改符号表。它改变了编译器外部状态，而不仅仅返回一个属性值。
\end{defbox}

\begin{defbox}[Attribute grammar]
属性文法是没有副作用的 SDD。其规则只用其他属性值和常量定义属性值。
\end{defbox}

\begin{defbox}[Annotated parse tree]
标注分析树是在 parse tree 的每个节点上标出属性值的树。它帮助我们观察属性如何沿树传播。
\end{defbox}

AST 也可以被标注成 decorated AST，例如为每个表达式标注类型，为每个变量标注符号表入口，为每个函数标注返回类型。

\section{依赖图与属性值计算顺序 \pg{29--33}}

\subsection{依赖图 \pg{29--30}}

\begin{defbox}[Dependency graph]
依赖图描述某棵 parse tree 中各属性实例之间的信息流。每个属性实例是图中的一个节点；若 \(X.a\) 的计算依赖 \(Y.b\)，则从 \(Y.b\) 指向 \(X.a\) 建一条有向边。
\end{defbox}

在计算某个属性前，必须先计算它依赖的所有属性。依赖图的作用就是确定这些属性值是否能被计算，以及以什么顺序计算。

例如声明：
\[
float\ a,b,c
\]
中，\(T.type=float\) 必须先算出，然后传给每层 \(L.inh\)，再用于对每个 id 执行 \(addtype(id.entry,L.inh)\)。

\subsection{拓扑序和环 \pg{31--33}}

合法属性求值顺序必须满足：若图中有边 \(M\to N\)，则 \(M\) 对应属性必须先于 \(N\) 计算。把有向图嵌入线性序列的过程称为拓扑排序（topological sort）。

\begin{notebox}[核心结论]
如果依赖图无环，则至少存在一个拓扑序；如果存在环，则没有拓扑序，也就没有办法在这棵 parse tree 上完成属性计算。
\end{notebox}

课件给出的循环例子：
\[
A\to B,\qquad A.s=B.i,\qquad B.i=A.s+1.
\]
这里 \(A.s\) 依赖 \(B.i\)，而 \(B.i\) 又依赖 \(A.s\)，形成环，无法求值。

综合属性一般需要先计算子节点再计算父节点；继承属性一般需要先计算父节点和依赖的兄弟节点，再计算当前节点。SDD 同时使用两类属性时，不保证自动存在合法顺序，因此需要 S-SDD 和 L-SDD 这样的受限子类。

\section{S-Attributed 与 L-Attributed Definitions \pg{34--37}}

\subsection{S-Attributed Definition \pg{35}}

\begin{defbox}[S-SDD]
若一个 SDD 中所有属性都是综合属性，则它是 S-attributed definition，简称 S-SDD。
\end{defbox}

S-SDD 的性质：

\begin{itemize}
  \item 属性都从子节点向父节点传递。
  \item 可以按任意自底向上顺序计算属性。
  \item 很适合在 bottom-up/LR parsing 中实现，因为规约时右部属性已在栈顶，正好可计算左部属性。
\end{itemize}

表达式求值文法就是典型 S-SDD。

\subsection{L-Attributed Definition \pg{36--37}}

\begin{defbox}[L-SDD]
L-attributed definition 允许综合属性和受限制的继承属性。对产生式
\[
A\to X_1X_2\cdots X_i\cdots X_n,
\]
\(X_i\) 的继承属性只能依赖：\(A\) 的继承属性、\(X_i\) 左侧符号 \(X_1,\ldots,X_{i-1}\) 的属性，以及 \(X_i\) 自身不成环的属性。
\end{defbox}

L-SDD 的直观含义是：产生式右部中信息只能从左向右传，不能依赖还没处理到的右兄弟。这样 top-down 解析时，走到某个符号前，它需要的继承属性已经可得。

示例：
\[
A\to BC,\qquad A.s=B.b,\qquad B.i=f(C.c,A.s).
\]
它不是 S-SDD，因为出现继承属性 \(B.i\)。它也不是 L-SDD，因为 \(B.i\) 依赖右兄弟 \(C.c\) 和左部综合属性 \(A.s\)，这些在处理 \(B\) 前通常还不可得。

\section{SDT 的实现与 S-SDD 的 LR 实现 \pg{38--47}}

\subsection{两类实现路径 \pg{39--40}}

SDT 可用两种方式实现：

\begin{enumerate}
  \item 先构造 parse tree 或 AST，再遍历树执行规则/动作。只要依赖无环，就可安排合适顺序，适用范围广。
  \item 在 parsing 过程中执行动作。此时动作顺序受分析器推导/规约顺序限制，只适合特定 SDD/SDT。
\end{enumerate}

课件强调两个重要可实现子类：

\begin{itemize}
  \item S-SDD，且底层文法可 LR 分析。
  \item L-SDD，且底层文法可 LL 分析。
\end{itemize}

\subsection{S-SDD 转 postfix SDT \pg{41--42}}

S-SDD 转 SDT 的规则很简单：把每个语义动作放在对应产生式的末尾。所有动作都在 RHS 末尾的 SDT 称为 postfix SDT。

例如：
\[
E\to E_1+T\quad \{E.val=E_1.val+T.val\}.
\]
在 LR 分析中，只有当 \(E_1+T\) 完整出现在栈顶时才会规约为 \(E\)，这时计算 \(E.val\) 所需属性都已经在栈中，所以动作可随规约执行。

\subsection{扩展 LR 分析栈 \pg{43--47}}

为了在 LR 规约时执行语义动作，栈记录不仅保存状态和文法符号，还要保存属性。对产生式：
\[
A\to XYZ,\qquad A.a=f(X.x,Y.y,Z.z),
\]
规约前栈顶有 \(X,Y,Z\) 及其属性；执行动作后，把这三项替换为 \(A\) 和 \(A.a\)。

对应栈操作可写成：
\[
stack[top-2].val=f(stack[top-2].val,stack[top-1].val,stack[top].val),
\]
\[
top=top-2.
\]

表达式 \(3*5+4\) 中，规约 \(T\to T*F\) 时，栈顶保存 \(T.val=3\)、\(*\)、\(F.val=5\)，于是计算 \(T.val=15\)；之后再与 \(+4\) 规约，得到最终表达式值 19。

\section{L-SDD 到 SDT、消除左递归与属性传递 \pg{48--60}}

\subsection{L-SDD 转 SDT 的规则 \pg{48--50}}

L-SDD 转成 SDT 时遵循：

\begin{enumerate}
  \item 若要计算某个 RHS 非终结符 \(A\) 的继承属性，把动作放在 \(A\) 出现之前。
  \item 若要计算 LHS 的综合属性，把动作放在产生式右部末尾。
\end{enumerate}

抽象形式：
\[
A\to B\ \{C.inh=\cdots\}\ C\ \{A.syn=\cdots\}.
\]
先准备 \(C\) 的继承属性，再进入 \(C\)；等 RHS 完成后，再计算 \(A\) 的综合属性。

例如：
\[
T\to F\ \{T'.inh=F.val\}\ T'\ \{T.val=T'.syn\}.
\]

\subsection{消除左递归时保留语义动作 \pg{51--55}}

Top-down 实现需要先消除左递归。若原文法：
\[
E\to E+T\mid E-T\mid T,
\]
带语义规则：
\[
E.val=E_1.val+T.val,\quad E.val=E_1.val-T.val.
\]
消左递归后：
\[
\begin{aligned}
E&\to T\ \{R.i=T.val\}\ R\ \{E.val=R.s\},\\
R&\to +T\ \{R_1.i=R.i+T.val\}\ R_1\ \{R.s=R_1.s\},\\
R&\to -T\ \{R_1.i=R.i-T.val\}\ R_1\ \{R.s=R_1.s\},\\
R&\to \eps\ \{R.s=R.i\}.
\end{aligned}
\]

这里 \(R.i\) 是“已经累计到目前为止的值”，\(R.s\) 是“把后续所有 \(+T/-T\) 都处理完后的值”。对输入 \(9-5+2\)，属性流为：
\[
R.i=9\to 4\to 6,\qquad E.val=6.
\]

一般规则：
\[
A\to A_1Y\ \{A.a=g(A_1.a,Y.y)\}\mid X\ \{A.a=f(X.x)\}
\]
可改为：
\[
A\to X\ \{R.i=f(X.x)\}\ R\ \{A.a=R.s\},
\]
\[
R\to Y\ \{R_1.i=g(R.i,Y.y)\}\ R_1\ \{R.s=R_1.s\}\mid \eps\ \{R.s=R.i\}.
\]

\section{L-SDD 在 LL/RDP/LR 中的实现 \pg{61--73}}

\subsection{递归下降实现 \pg{61--62}}

递归下降 parser 通常为每个非终结符写一个函数。L-SDD 在这种实现中很自然：

\begin{itemize}
  \item 继承属性作为函数参数传入。
  \item 综合属性作为函数返回值返回。
  \item 中间语义动作使用局部变量保存。
\end{itemize}

例如 \(T'\) 的函数可写成接受 \(T'.inh\) 参数，返回 \(T'.syn\)。若当前 token 是 \(*\)，就读取 \(F.val\)，计算 \(T_1'.inh=T'.inh*F.val\)，递归调用 \(T_1'\)，最后返回 \(T_1'.syn\)。若当前 token 是结束符或 Follow 符号，就取 \(\eps\) 分支并返回 \(T'.inh\)。

\subsection{非递归 LL 预测分析实现 \pg{63--67}}

表驱动 LL parser 用栈模拟最左推导。为了执行 L-SDD，需要扩展分析栈：

\begin{itemize}
  \item Action-record：表示待执行动作。
  \item Synthesize-record：保存非终结符的综合属性。
  \item 非终结符的继承属性通常放在该符号对应的栈记录中。
  \item 非终结符的综合属性可放在紧邻该符号的单独记录中。
\end{itemize}

难点是：当栈顶非终结符 \(A\) 展开为 \(BC\) 时，\(A\) 的记录会被弹出；如果后面计算 \(C.i\) 还需要 \(A.i\) 或 \(B\) 的属性，就必须提前把这些值复制到后续动作记录中。

因此课件总结两条操作原则：

\begin{enumerate}
  \item 左部展开时，若后续动作需要左部继承属性，要把它复制给后面的动作记录。
  \item 综合属性出栈时，若后续动作需要它，也要复制给后面的动作记录。
\end{enumerate}

\subsection{L-SDD 在 LR 中的 marker 技巧 \pg{68--73}}

LR 比 LL 语法分析能力更强；S-SDD 在 LR 中很好实现，因为动作都在产生式末尾。但 L-SDD 的动作可能嵌在 RHS 中间，这不符合 LR 规约时才执行动作的自然机制。

问题例子：
\[
A\to X\alpha\ \{Y.in=X.s\}\ Y\mid X\beta\ \{Y.in=X.s\}\ Y,
\]
若 \(|\alpha|\ne|\beta|\)，则 \(X.s\) 相对 \(Y\) 的栈位置不同，无法用固定偏移统一取得。

解决方案是引入 marker 非终结符：
\[
A\to X\alpha M_1Y\mid X\beta M_2Y,
\]
\[
M_1\to\eps\ \{M_1.s=stack[top-|\alpha|].s\},\quad
M_2\to\eps\ \{M_2.s=stack[top-|\beta|].s\}.
\]
Marker 是只产生 \(\eps\) 的占位非终结符。它把“中间动作”转换为“某个空产生式规约时执行的末尾动作”，从而重新适配 LR 的规约机制。

抽象改写：
\[
A\to \{B.i=f(A.i)\}\ BC
\]
变为：
\[
A\to MBC,\qquad M\to\eps\ \{M.i=A.i;\ M.s=f(M.i)\}.
\]

\section{语义分析应用：绑定、作用域、符号表 \pg{74--93}}

\subsection{Binding 与 Scope \pg{76--80}}

\begin{defbox}[Binding]
绑定是把标识符使用点匹配到其定义的过程。定义通常把标识符与内存位置、类型、函数体等信息关联起来。
\end{defbox}

若程序中有多个同名定义，编译器必须判断使用点绑定到哪一个。例子中外层有 \(char\ x\)，内层块又有 \(int\ x\)，在不同位置出现的 \(x=x+1\) 可能对应不同变量。

\begin{defbox}[Scope]
作用域是某个定义可用于绑定的程序区域。在该区域内，标识符使用会绑定到该定义；若使用点不在任何定义作用域内，则产生未定义错误。
\end{defbox}

两类作用域：

\begin{longtable}{L{0.18\linewidth}L{0.36\linewidth}L{0.36\linewidth}}
\toprule
类型 & 决定方式 & 特点 \\
\midrule
静态作用域 / lexical scoping & 由程序文本结构决定，绑定到最近的外层定义 & 编译期可分析，错误少，代码更高效；C/C++、Java、Python、JavaScript 常见 \\
动态作用域 & 由运行时执行路径决定，绑定到当前执行中最近定义 & 运行时才清楚，人和编译器都较难判断；Bash、早期 Lisp 等出现过 \\
\bottomrule
\end{longtable}

多数语言后来倾向静态作用域，因为绑定可从代码结构直接看出，编译器也可提前优化。

\subsection{符号表是什么 \pg{81--82}}

\begin{defbox}[Symbol table]
符号表是记录程序中所有符号信息的数据结构。每个条目代表一个标识符定义，维护当前程序点可达的定义，用于把标识符使用绑定到正确定义。
\end{defbox}

常见操作：

\begin{itemize}
  \item \texttt{enter\_scope()}：进入新作用域。
  \item \texttt{exit\_scope()}：退出当前作用域。
  \item \texttt{find\_symbol(x)}：查找 \(x\) 的当前可见定义。
  \item \texttt{add\_symbol(x)}：把 \(x\) 的定义加入符号表。
  \item \texttt{check\_symbol(x)}：检查 \(x\) 是否已在当前作用域定义。
\end{itemize}

符号表通常在生成可执行文件后可丢弃，因为机器码不再使用变量名；但为了调试器显示源码变量名，编译时可保留 debug symbol table，例如 GCC 的 \texttt{-g}。

\subsection{符号表数据结构 \pg{83--87}}

符号表访问很频繁，会影响编译前端性能。常见结构对比：

\begin{longtable}{L{0.18\linewidth}L{0.26\linewidth}L{0.26\linewidth}L{0.20\linewidth}}
\toprule
结构 & 空间 & 插入/删除 & 查找 \\
\midrule
数组 & 几乎无额外空间 & \(O(n)\) & \(O(n)\) \\
链表 & 额外指针空间 & \(O(1)\) & \(O(n)\)，可把最近访问项移到表头优化 \\
平衡二叉树 & 比数组/链表更多空间 & \(O(\log n)\) & \(O(\log n)\)，最坏可退化 \\
哈希表 & 表大小 \(M\) 可能大于符号数 \(N\) & 平均 \(O(1)\) & 平均 \(O(1)\)，冲突链过长会退化 \\
\bottomrule
\end{longtable}

多数编译器选择哈希表，因为符号查找非常频繁，平均 \(O(1)\) 的访问时间更有价值。哈希冲突通常用链式法解决。

\subsection{处理多作用域 \pg{88--90}}

概念上，每个作用域都需要一张独立符号表。活动作用域可用作用域栈管理：

\begin{itemize}
  \item 进入作用域时创建新符号表，并把指针压入栈。
  \item 退出作用域时弹出该符号表。
  \item 查找标识符时从栈顶向下搜索，先找到的就是最近内层定义。
\end{itemize}

这种结构自然匹配嵌套块结构，例如 \(\{C_1\{C_2\{C_3\}\}\{C_4\}\}\)。栈顶总是最内层作用域。

\subsection{符号表条目和类型信息 \pg{91--93}}

符号表条目通常包含：

\begin{itemize}
  \item name：标识符字符串。
  \item kind：变量、函数、结构体类型、类类型等。
  \item attributes：随 kind 不同而变化的属性列表。
\end{itemize}

不同符号的属性不同：

\begin{itemize}
  \item 变量：类型、地址。
  \item 函数：函数原型、函数体地址。
  \item 结构体类型：字段名、字段类型、字段大小。
  \item 类类型：类作用域的符号表。
\end{itemize}

类型信息可很复杂。例如数组可能有多维大小，结构体可能包含多个字段。因此常把类型信息放在属性列表中。

类型信息用于：

\begin{itemize}
  \item 代码生成：变量赋值是 4 字节拷贝还是 8 字节拷贝；加法是整数加法还是浮点加法。
  \item 函数调用翻译：参数和返回值类型决定调用约定。
  \item 代码优化：帮助编译器理解程序语义。
  \item 类型检查：判断程序类型使用是否合法。
\end{itemize}

\section{类型、类型检查与类型系统 \pg{94--103}}

\subsection{Type 与 Type Checking \pg{94}}

\begin{defbox}[Type]
类型是一组值以及定义在这些值上的一组操作。例如 int 包含整数值并支持整数加减乘除；char* 支持指针相关操作和数组下标访问。
\end{defbox}

\begin{defbox}[Type checking]
类型检查验证程序中操作和操作数类型是否一致。若所有操作都适用于其操作数类型，则程序在类型意义上是一致的。
\end{defbox}

例子：

\begin{itemize}
  \item \texttt{char *str = "Hello";} 中，\texttt{str[2]} 合法，因为 char* 可做下标访问；\texttt{str/2} 不合法，因为指针不支持除法。
  \item \texttt{int pi = 3;} 中，\texttt{pi/2} 合法；\texttt{pi=3.14} 类型不一致，除非语言允许隐式类型转换。
\end{itemize}

\subsection{静态与动态类型检查 \pg{95--96}}

\begin{defbox}[静态类型检查]
静态类型检查在编译时通过代码分析判断类型一致性。它依赖声明和符号表收集的信息。
\end{defbox}

例如：
\[
int\ a,b,c;\quad a=b+c;
\]
若 \(b,c\) 都是 int，编译器可证明 \(b+c\) 仍是 int，赋给 \(a\) 合法。

静态检查不能发现所有错误，例如整数溢出、某些强制类型转换后的实际风险。

\begin{defbox}[动态类型检查]
动态类型检查在运行时检查值携带的类型标签。每次操作前先检查 type tag，合法才执行，否则触发运行时类型错误。
\end{defbox}

动态检查可处理编译时无法确定的情况，例如数组越界、运行时 downcasting 等，但会带来运行时开销。总体上，能在编译时发现的错误越多越好。

\subsection{Static Typing 与 Dynamic Typing \pg{97--99}}

静态/动态 typing 是语言类型系统特征；静态/动态 type checking 是执行检查的方法。

\begin{longtable}{L{0.18\linewidth}L{0.36\linewidth}L{0.36\linewidth}}
\toprule
类型系统 & 特点 & 优缺点 \\
\midrule
Static typing & 变量有静态类型，只能保存该类型的值；C/C++、Java 常见 & 程序员约束更多，但 bug 更少、运行更快 \\
Dynamic typing & 变量可在运行时保存不同类型值；Python、JavaScript、PHP 常见 & 代码更灵活更短，适合脚本和原型；但错误可能延迟到运行时，执行开销更大 \\
\bottomrule
\end{longtable}

静态检查不只用于静态类型语言：动态类型语言若编译器能推断类型，也可做静态检查。动态检查也不只用于动态类型语言：静态类型语言中的 downcasting、数组越界等也可能需要运行时检查。

\section{课件总结 \pg{100--103}}

本讲知识点可归纳为三组：

\begin{enumerate}
  \item SDD/SDT 基础：属性、综合属性、继承属性、依赖图、拓扑序、S-SDD、L-SDD。
  \item SDD/SDT 实现：S-SDD 可转 postfix SDT 并在 LR 规约时执行；L-SDD 可在递归下降中用参数/返回值，在 LL 中扩展栈，在 LR 中引入 marker。
  \item 语义分析应用：绑定、作用域、符号表、类型信息、类型检查、静态/动态类型系统。
\end{enumerate}

推荐阅读对应 Dragon Book 第 5 章 5.1--5.5 节，分别覆盖语法制导翻译介绍和 top-down/bottom-up LR 中翻译方案的实现。

\section{练习题与解答 \pg{104--129}}

\subsection*{题 1：无二义文法的最左/最右推导分析树是否一样？}

\textbf{答案：} 一样。无二义文法中，一个句子只有一棵分析树。最左推导和最右推导只是推导顺序不同，分析树会隐去具体推导顺序。

\subsection*{题 2：如何判断 LL(1) 文法？}

对每个非终结符的候选式 \(A\to\alpha\mid\beta\)，要求：
\[
\FIRST(\alpha)\cap\FIRST(\beta)=\emptyset.
\]
若 \(\beta\Rightarrow^*\eps\)，还要求：
\[
\FIRST(\alpha)\cap\FOLLOW(A)=\emptyset.
\]
也就是说，一个 lookahead 符号不能同时指向两条候选式。

\subsection*{题 3：LL(1) 表行列数由什么决定？}

\textbf{答案：} 行数由非终结符个数决定；列数由终结符个数加 1 决定，多出的 1 是结束符号 \(\$\)。

\subsection*{题 4：活前缀}

设最右推导：
\[
S\Rightarrow aBb\Rightarrow aaAb
\]
句柄为 \(aA\)。活前缀是不越过句柄右端的前缀，因此为：
\[
a,\quad aa,\quad aaA.
\]

\subsection*{题 5：短语、直接短语与句柄}

文法：
\[
E\to T\mid E+T,\quad T\to F\mid T*F,\quad F\to (E)\mid a\mid b\mid c.
\]
句型 \(a*b+c\) 的结果：

\[
\text{短语： } a*b+c,\ a*b,\ a,\ b,\ c.
\]
\[
\text{直接短语： } a,\ b,\ c.
\]
\[
\text{句柄： } a.
\]

\subsection*{题 6：构造 LL(1) 分析表}

文法：
\[
S\to A,\qquad A\to\eps,\qquad A\to bbA.
\]
\[
\FIRST(A)=\{b,\eps\},\qquad \FOLLOW(S)=\FOLLOW(A)=\{\$\}.
\]

\begin{center}
\begin{tabular}{c|cc}
\toprule
 & \(b\) & \(\$\) \\
\midrule
\(S\) & \(S\to A\) & \(S\to A\) \\
\(A\) & \(A\to bbA\) & \(A\to\eps\) \\
\bottomrule
\end{tabular}
\end{center}

\subsection*{题 7：文法 \texorpdfstring{\(S\to A,\ A\to A+A\mid B++,\ B\to y\)}{S->A} 的 LR 表与 y++}

课件给出的关键分析表中，状态 5 在输入 \(+\) 时有 \(r2/s4\) 冲突，因此该文法不是 LR(1)。可通过改写文法、声明结合性/优先级或指定冲突解决策略处理。

对句子 \(y++\)：

\begin{center}
\small
\begin{tabular}{llll}
\toprule
状态栈 & 符号栈 & 输入 & 动作 \\
\midrule
0 & \(\$\) & \(y++\$\) & \(s1\) \\
0 1 & \(\$y\) & \(++\$\) & \(r4: B\to y\) \\
0 & \(\$B\) & \(++\$\) & \(\GOTO(0,B)=3\) \\
0 3 & \(\$B\) & \(++\$\) & \(s6\) \\
0 3 6 & \(\$B+\) & \(+\$\) & \(s7\) \\
0 3 6 7 & \(\$B++\) & \(\$\) & \(r3: A\to B++\) \\
0 & \(\$A\) & \(\$\) & \(\GOTO(0,A)=2\) \\
0 2 & \(\$A\) & \(\$\) & acc \\
\bottomrule
\end{tabular}
\end{center}

因此 \(y++\) 是该文法的一个句子。

\subsection*{题 8：文法 \texorpdfstring{\(S\to Sb\mid bAa,\ A\to aSc\mid aSb\mid a\)}{S grammar} 的 LR/SLR 分析}

增广文法：
\[
(0)\ S'\to S,\quad
(1)\ S\to Sb,\quad
(2)\ S\to bAa,\quad
(3)\ A\to aSc,\quad
(4)\ A\to aSb,\quad
(5)\ A\to a.
\]

LR(0) 项集可写为：

\begin{small}
\begin{longtable}{L{0.12\linewidth}L{0.78\linewidth}}
\toprule
状态 & 项 \\
\midrule
\(I_0\) & \(S'\to\cdot S;\ S\to\cdot Sb;\ S\to\cdot bAa\) \\
\(I_1\) & \(S\to b\cdot Aa;\ A\to\cdot aSc;\ A\to\cdot aSb;\ A\to\cdot a\) \\
\(I_2\) & \(S'\to S\cdot;\ S\to S\cdot b\) \\
\(I_3\) & \(S\to\cdot Sb;\ S\to\cdot bAa;\ A\to a\cdot Sc;\ A\to a\cdot Sb;\ A\to a\cdot\) \\
\(I_4\) & \(S\to bA\cdot a\) \\
\(I_5\) & \(S\to Sb\cdot\) \\
\(I_6\) & \(S\to S\cdot b;\ A\to aS\cdot c;\ A\to aS\cdot b\) \\
\(I_7\) & \(S\to bAa\cdot\) \\
\(I_8\) & \(S\to Sb\cdot;\ A\to aSb\cdot\) \\
\(I_9\) & \(A\to aSc\cdot\) \\
\bottomrule
\end{longtable}
\end{small}

该文法不是 LR(0)。原因是 LR(0) 完成项会对所有终结符规约，因而产生冲突：课件标出有移进-规约冲突和归约-归约冲突。按上面的编号，\(I_3\) 有 \(A\to a\cdot\) 且还能在 \(b\) 上移进；\(I_8\) 同时有两个完成项。

该文法是 LR(1)，因为它是 SLR(1)，自然也属于 LR(1)。关键 FOLLOW 集：
\[
\FOLLOW(S)=\{b,c,\$\},\qquad \FOLLOW(A)=\{a\}.
\]
SLR(1) 表如下：

\begin{center}
\small
\begin{tabular}{c|cccc|cc}
\toprule
State & \(a\) & \(b\) & \(c\) & \(\$\) & \(S\) & \(A\) \\
\midrule
0 &  & \(s1\) &  &  & 2 &  \\
1 & \(s3\) &  &  &  &  & 4 \\
2 &  & \(s5\) &  & acc &  &  \\
3 & \(r5\) & \(s1\) &  &  & 6 &  \\
4 & \(s7\) &  &  &  &  &  \\
5 &  & \(r1\) & \(r1\) & \(r1\) &  &  \\
6 &  & \(s8\) & \(s9\) &  &  &  \\
7 &  & \(r2\) & \(r2\) & \(r2\) &  &  \\
8 & \(r4\) & \(r1\) & \(r1\) & \(r1\) &  &  \\
9 & \(r3\) &  &  &  &  &  \\
\bottomrule
\end{tabular}
\end{center}

表中没有同一格多动作，因此是 SLR(1)。

\subsection*{题 9：L-属性定义是否每个属性都是继承属性？}

\textbf{答案：} 错。L-属性定义允许综合属性，也允许满足从左到右依赖限制的继承属性，并不是所有属性都是继承属性。

\subsection*{题 10：RHS 中间有语义动作是否只能用递归下降实现？}

\textbf{答案：} 错。中间动作确实更适合 top-down 直接执行，但不代表无法用 LR。可以通过 marker 非终结符把嵌入动作转化为空产生式上的 postfix 动作，再用 LR 栈操作实现。

\subsection*{题 11：SDD 输出题 \texorpdfstring{\((a,(a))\)}{(a,(a))}}

文法：
\[
S'\to S,\quad S\to (L)\mid a,\quad L\to L_1,S\mid S.
\]
语义动作：
\[
S'\to S\ \{print(S.num)\},\quad
S\to (L)\ \{S.num=L.num+1\},\quad
S\to a\ \{S.num=0\},
\]
\[
L\to L_1,S\ \{L.num=L_1.num+S.num\},\quad
L\to S\ \{L.num=S.num\}.
\]

输入 \((a,(a))\) 中，左侧 \(a\) 的 \(S.num=0\)，右侧 \((a)\) 的内部 \(a\) 为 0，因此 \((a)\) 对应的 \(S.num=1\)。外层列表的 \(L.num=0+1=1\)，外层括号使最终 \(S.num=L.num+1=2\)。输出为：
\[
2.
\]

\subsection*{题 12：判断属性规则是否满足 S-SDD / L-SDD}

产生式 \(A\to BCD\)，四个非终结符都有综合属性 \(s\) 和继承属性 \(i\)。

\begin{longtable}{L{0.31\linewidth}L{0.26\linewidth}L{0.33\linewidth}}
\toprule
规则 & S-SDD? & L-SDD? \\
\midrule
\(A.s=B.s+D.s\) & 满足，只有综合属性 & 满足，没有继承属性冲突 \\
\(A.s=B.i+C.s\) & 不满足，使用继承属性 \(B.i\) & 满足，未违反从左到右限制 \\
\(A.s=D.i,\ B.i=A.s+C.s,\ C.i=B.s,\ D.i=B.i+C.i\) & 不满足，存在继承属性 & 不满足，\(B.i\) 依赖右兄弟 \(C.s\)，且依赖左部综合属性 \(A.s\) \\
\bottomrule
\end{longtable}

\subsection*{题 13：非标准二进制数值系统}

文法：
\[
S\to T\mid\eps,\qquad T\to 0T\mid 1T\mid 0\mid 1.
\]
从右至左，非零二进制位对应的权值符号交替变化。只使用 \(val\) 和 \(tmp\) 两个属性，可按课件表填写：

\begin{center}
\small
\begin{tabular}{ll}
\toprule
产生式 & 语义动作 \\
\midrule
\(S\to T\) & \(S.val=T.val\) \\
\(S\to\eps\) & \(S.val=0\) \\
\(T_1\to 0T_2\) & \(T_1.tmp=2*T_2.tmp;\quad T_1.val=T_2.val\) \\
\(T_1\to 1T_2\) & \(T_1.tmp=-2*T_2.tmp;\quad T_1.val=T_2.val+T_1.tmp\) \\
\(T\to 0\) & \(T.tmp=-1;\quad T.val=0\) \\
\(T\to 1\) & \(T.tmp=1;\quad T.val=1\) \\
\bottomrule
\end{tabular}
\end{center}

直观解释：\(tmp\) 表示当前位置对应的带符号权值；从右向左移动一位时，权值绝对值乘 2，符号交替。若当前位是 0，只传递后续值；若当前位是 1，把当前带符号权值加入 \(val\)。

\section*{复习建议}
\addcontentsline{toc}{section}{复习建议}

复习本讲建议抓住四组对应关系：

\begin{itemize}
  \item SDD vs SDT：一个是规则说明，一个是带执行位置的方案。
  \item 综合属性 vs 继承属性：一个向上传，一个向下/横向传。
  \item S-SDD vs L-SDD：一个只允许综合属性，一个允许受限继承属性。
  \item 静态 vs 动态：作用域、类型系统、类型检查都可能有静态/动态之分，但含义不同。
\end{itemize}

做题时优先画出属性依赖关系：若依赖图无环，再看是否满足 S-SDD 或 L-SDD 的限制；若要和 parser 同步执行，再考虑动作应该放在产生式末尾、中间，还是需要 marker 改写。

\end{document}
